\section{Условие}
Создать базу знаний <<Собственники>>, дополнив (и минимально изменив) базу знаний, хранящую знания (лаб. 7):

\begin{itemize}
    \item <<Телефонный справочник>>: Фамилия, № тел, Адрес --- структура (Город, Улица, № дома, № кв);
    \item <<Автомобили>>: Фамилия владельца, Марка, Цвет, Стоимость и др.;
    \item <<Вкладчики банков>>: Фамилия, Банк, Счет, Сумма и др.
\end{itemize}

Преобразовать знания об автомобиле к форме знаний о собственности. 
Дополнить базу информацией о других видах собственности:
\begin{itemize}
    \item Строение: вид, стоимость и другие его характеристики;
    \item Участок: вид, стоимость и другие его характеристики;
    \item Водный транспорт: стоимость и другие его характеристики.
\end{itemize}

Владелец может иметь, но только один объект каждого вида собственности 
(это касается автомобиля, строения, участка и водного транспорта).

Используя конъюнктивное правило и разные формы задания одного вопроса, обеспечить возможность поиска:
\begin{enumerate}
    \item Названий всех объектов собственности заданного субъекта;
    \item Названий и стоимости всех объектов собственности заданного субъекта;
    \item Разработать правило, позволяющее найти суммарную стоимость всех объектов собственности заданного субъекта.
\end{enumerate}